\section{Synthetic study: where the adjudication goes}
\label{sec:synthetic}

This study supplies the evidence for Findings~1 and~3, and the ablation behind
Finding~2; the findings themselves are stated in Section~\ref{sec:findings}.

\subsection{Definitions}
\label{sec:metrics}

Confusions are defined over the \emph{verdict}: a \emph{false proof} is the
strict subset of false alarms promoted to \texttt{violation-proven}---an error
the tool asserts as established. We also report the \emph{detection-keep rate}
(share of fired true violations still alarmed) and count only abstentions
justified by an unreadable oracle.

\subsection{What the constructed system shows}

\paragraph{Attribution is clean} On the anchored configuration the false-proof
rate is \ResFpEthreeHits{}/\ResFpEthreeTotal{}. With that denominator the rate is
unstable by construction, so we attach to it no interval---nor a percentage,
which would travel out of
context. With the legitimate-sharing
exception absent the rate is \ResBaseFpRate{}
(\ResBaseFpHits{}/\ResBaseFpTotal{}), and with the channel listed it is
\ResListedCtlFpRate{} (\ResListedCtlFpHits{}/\ResListedCtlFpTotal{}). The
phenomenon is therefore attributable to a legitimate sharing exception the MR
cannot know about, not to execution itself. A $2{\times}2$ grid over channel
visibility and enforcement correctness tells the same story from the other
side: on unlisted cells the false-proof rate is
\ResBEThreevulnZerounlistedFpPct{}\% with no false negatives; on listed cells it
drops to \ResBEThreevulnZerolistedFpPct{}\% but one of the five violating records
in that cell is a false negative (\ResBEThreevulnOnelistedTp{} detected,
\ResBEThreevulnOnelistedFn{} missed). Avoiding false proofs and keeping detection are
controlled by one switch, and the switch cannot be thrown both ways at once.

\paragraph{Filtering exceptions does not help}
The exception ladder---raw, one exception level, two---leaves the anchored rate
at \ResLadderrawRate{}, \ResLaddereOneRate{}, \ResLaddereOneeTwoRate{}: it does
not move. The exception levels we can express are themselves functions of the
same reachability model, so they inherit its blind spot.

\paragraph{The framework cannot see the channel that matters}
When a legitimate share was created through the framework's facility, its log
holds \ResObsFwLogRecords{} records and v2 produces \ResObsFwVtwoFalseProof{}
false proofs (rate \ResObsFwVtwoFpRate{}); when the same share was created out
of band, the log holds \ResObsOobLogRecords{} records---no evidence the share
exists---and v2's false proofs rise to \ResObsOobVtwoFalseProof{} of
\ResObsOobFalseAlarm{} alarms (rate \ResObsOobVtwoFpRate{}), every one promoted
to \texttt{violation-proven}. The MR cannot distinguish ``this share does not
exist'' from ``this share exists but was not crawled''; its only representation
of the first is the second. On S7, where a genuine violation \emph{does} sit
behind the out-of-band share, the same blindness makes the alarm right for the
wrong reason (\ResObsOobBolaVtwoProvenBola{} of \ResObsOobBolaBolaTotal{}
reached)---and that coincidence of alarms with violations is what an
input-level metric cannot unpack. The same tension appears within a single
MR: for S3 (must not alarm) vs S5 (must alarm), no configuration in the family
is correct on both---the choice of switch is a choice of \emph{which} error to
make (Fig.~\ref{fig:scenarios}).

\begin{figure*}[t]
\centering
\includegraphics[width=0.78\textwidth]{fig2_scenarios.pdf}
\caption{The errors are scenario-structured. The bookkeeping oracle's false proofs
concentrate on~S6, the legitimate share created out of band and so invisible to
the framework's own log, while the state-calibrated v3 abstains in two places:
where its surface is not readable, and where a readable surface shows a covering
grant it declines to override. Both panels are recomputed from the same records
as Table~\ref{tab:ablation}.}
\label{fig:scenarios}
\end{figure*}

\paragraph{Stratification, not aggregation}
Aggregated, v3 looks best (recovery \ResVdVThreeRecoveryRate{} against
\texttt{degrade\_all}'s keep rate \ResVdDegradeallKeepRate{}); the reading is
wrong, because v3 changes regime with oracle readability. Stratified over the
\ResIntroOnN{}-record readable grid and its \ResIntroOffN{}-record absent grid (a
representative grid; Section~\ref{sec:stratum} re-runs the same comparison over
all \ResSynRecords{} records): on the readable grid it removes v2's false proofs
(\ResIntroOnVthreeFpN{} vs.~\ResIntroOnVtwoFpN{}) while still alarming on all
\ResIntroOnVthreeProvenBola{} fired true violations; on the absent grid
(\ResIntroOffN{}) it alarms on none of the \ResIntroOffVtwoProvenBola{} true
violations v2 still alarms on, abstaining \ResIntroOffAbstain{} times. Pooling
converts honest abstention into apparent performance, so v3's numbers may appear
only stratified.

\subsection{Fair ablation: does the metamorphic step add adjudication information?}
\label{sec:synthetic-ablation}

Our earlier iteration proposed v3 as a mechanism, resting on the claim that the
metamorphic step was necessary. We tested that claim by removing the step and
substituting the most naive consumer of the same observable. The distinction
matters: \texttt{direct\_only}---oracle yes, access outcome
ignored---produces \ResAbDirectOnlyFP{} false proofs against v3's
\ResAbVThreeFP{} (precision \ResAbDirectOnlyPrec{}), and rejecting it proves
nothing because it discards half the evidence. The fair control is
\texttt{direct\_observe}: a succeeding ordinary request plus the same readable
surface, consuming exactly what the MR-based adjudicator consumes.

\begin{table*}[t]
\centering
\caption{End-to-end confusion matrix over all \ResSynRecords{} records
(\ResAbNBola{} true violations, \ResAbNNormal{} correct cases). The records are
a fixed enumeration, not a sample (Section~\ref{sec:protocol}, statistical
treatment), so no inferential
interval is attached; where a denominator is small we say so rather than
interval it. \texttt{direct\_only} is a deliberately unfair control, shown to
make that point.}
\label{tab:ablation}
\begin{tabular}{lrrrrrrl}
\toprule
method & TP & FP & FN & TN & precision & recall & F1 \\
\midrule
raw (no exception handling) & \ResAbRawTP{} & \ResAbRawFP{} & \ResAbRawFN{} & \ResAbRawTN{} & \ResAbRawPrec{} & \ResAbRawRec{} & \ResAbRawFone{} \\
v1 (existence filter) & \ResAbVOneTP{} & \ResAbVOneFP{} & \ResAbVOneFN{} & \ResAbVOneTN{} & \ResAbVOnePrec{} & \ResAbVOneRec{} & \ResAbVOneFone{} \\
v2 (existence + coverage) & \ResAbVTwoTP{} & \ResAbVTwoFP{} & \ResAbVTwoFN{} & \ResAbVTwoTN{} & \ResAbVTwoPrec{} & \ResAbVTwoRec{} & \ResAbVTwoFone{} \\
v3 (state-calibrated) & \ResAbVThreeTP{} & \ResAbVThreeFP{} & \ResAbVThreeFN{} & \ResAbVThreeTN{} & \ResAbVThreePrec{} & \ResAbVThreeRec{} & \ResAbVThreeFone{} \\
\texttt{degrade\_all} & \ResAbDegradeAllTP{} & \ResAbDegradeAllFP{} & \ResAbDegradeAllFN{} & \ResAbDegradeAllTN{} & \ResAbDegradeAllPrec{} & \ResAbDegradeAllRec{} & \ResAbDegradeAllFone{} \\
\texttt{direct\_only} & \ResAbDirectOnlyTP{} & \ResAbDirectOnlyFP{} & \ResAbDirectOnlyFN{} & \ResAbDirectOnlyTN{} & \ResAbDirectOnlyPrec{} & \ResAbDirectOnlyRec{} & \ResAbDirectOnlyFone{} \\
\texttt{direct\_observe} & \ResAbDirectObserveTP{} & \ResAbDirectObserveFP{} & \ResAbDirectObserveFN{} & \ResAbDirectObserveTN{} & \ResAbDirectObservePrec{} & \ResAbDirectObserveRec{} & \ResAbDirectObserveFone{} \\
\bottomrule
\end{tabular}
\end{table*}

Table~\ref{tab:ablation} concludes the ablation. The first three rows are a
ladder, and they must be read in order: with no exception handling at all,
\texttt{raw} yields \ResAbRawFP{} false proofs; the existence filter v1 cuts
them to \ResAbVOneFP{}, but pays for it by abstaining on true violations (FN
rises to \ResAbVOneFN{}); adding operation coverage (v2) restores the true
violations (FN \ResAbVTwoFN{}) at the same false-proof count. Neither filter is
therefore a free improvement, and v3's zero in the FP column is bought with the
FN column. \texttt{direct\_observe} has
\ResAbDirectObserveFP{} false proofs against v3's \ResAbVThreeFP{}---a tie at
zero, not a win. The comparison is paired, so we report the discordance rather
than an inferential interval: over the \ResAbNBola{} true-violation records the
two methods disagree on \ResAbDisagreementN{}, and on every one of them
\texttt{direct\_observe} is right---it detects a violation that v3 misses or
abstains on, its reachability precondition blocking a detection the plain
request makes. The aggregate recalls \ResAbDirectObserveRec{} and
\ResAbVThreeRec{} restate that asymmetry rather than measure a difference, and
no margin was pre-specified, so no formal hypothesis-test claim attaches. We
found no incremental adjudication advantage of the metamorphic step over direct
observation in these oracle-available scenarios, and that absence provides no
support for an incremental adjudication benefit in these scenarios:
the step would have to be better at something the plain request cannot do.
Selective presentation would have hidden the cost: v3's precision of
\ResAbVThreePrec{} was bought with \ResAbVThreeFN{} false negatives (F1
\ResAbVThreeFone{}), an exchange of \ResAbVOneFP{} false alarms for missed
violations dressed as an improvement.

\subsection{Coverage, selective risk, and what a pooled precision does not say}
\label{sec:coverage}

Table~\ref{tab:ablation} scores the alarm. A method that abstains is counted as
not having alarmed, so the abstention lands in TN or in FN: for the question
``does the tool report a violation'' that is the right scoring, but it is not the
scoring for ``does the method decide''. Table~\ref{tab:coverage} rescores the
same \CvgN{} records on the decision and reports coverage, abstention rate,
selective risk and request cost.

v3 decides \CvgVthreeDecided{} of \CvgN{} records---coverage
\CvgVthreeCoveragePct{}\%---and abstains on \CvgVthreeAbstain{},
\CvgVthreeAbstBola{} of them on true violations and \CvgVthreeAbstNormal{} on
records with none. Its precision of $1.00$ and its selective risk of
\CvgVthreeRiskPct{}\% are therefore properties of a \CvgVthreeCoveragePct{}\% coverage: over the
remaining two records in five the method claims nothing, and the pooled precision
is not an error rate over all records.

The abstentions are not all of one kind (Table~\ref{tab:coverage}, panel b).
\CvgVthreeUnread{} are structural---the surface returns nothing, and there is
nothing for the method to consult. The other \CvgVthreeCover{} fall on a
\emph{readable} surface that shows a covering grant: the access is explained, and
the method still declines to assert its absence. A conservative method may prefer
silence to a positive ``no violation'', and we do not score that as an error; but
it is a non-decision, and it cannot be read as a correct rejection. The
comparison with v2 makes the arithmetic visible. v2 promotes \CvgVtwoFP{} false
proofs; v3 declines \CvgVtwoFPUnread{} of them because the surface is unreadable
and \CvgVtwoFPCover{} because it is readable and shows the grant. Only the first
\CvgVtwoFPUnread{} are the availability gate the method is named for. The rest is
a false proof converted into a non-decision.

The other side of the ledger is dearer. \CvgVthreeAbstBola{} abstentions fall on
true violations, and with the \CvgVthreeMissed{} records on which no alarm is
raised at all, v3 leaves \CvgVthreeFNunproven{} of \CvgViolations{} violations
unproven---the count the review derived from the pooled table.

\texttt{direct\_only} is scored under its own adjudication rule, which is also
the rule Table~\ref{tab:ablation} uses: assert a violation whenever the
authorization surface reports no granting channel, regardless of the request
outcome. On the \CvgN{} records it decides \CvgDirOnlyDecided{} (the same
abstention set as \texttt{direct\_observe}, so the two differ in rule, not in
coverage) and it is right on \CvgDirOnlyTP{} true violations. Its
\CvgDirOnlyFP{} false proofs are not, however, strict false proofs in the sense
defined in Section~\ref{sec:protocol}: every one of them falls on a request the
server in fact denied, so no unauthorized access resulted
(\CvgDirOnlyFpNoEffect{} of \CvgDirOnlyFP{} without violation effect,
\CvgDirOnlyFpEffect{} on allowed requests). They are invalid violation
claims---the oracle's ``no granting channel'' read as a violation claim that the
request outcome contradicts---and its selective risk over the decided set is
\CvgDirOnlyRiskPct\%, not zero. Keeping the rule fixed across both tables is what
exposes this cost; a table that silently added the request-success condition to
this row would have reported a zero risk that the method never earned.
The overlap of v3's decisions with \texttt{direct\_only}'s covers
\CvgJointVthreeDirect{} records---the \CvgObsDecided{} records on which the direct
method decides minus the \CvgVthreeCover{} records on which v3 additionally
declines on a readable covering surface---and the violation-detection difference
on comparable denominators is stated with its own denominator, not with the
\CvgObsDecided{} oracle-readable count.

\texttt{direct\_observe} is scored on the same contract, and that changes what
its earlier numbers meant. It abstains when the probe is unreadable, exactly as
the direct adjudicator does, so it decides \CvgObsDecided{} records---coverage
\CvgObsCoveragePct{}\%---and abstains on \CvgObsAbstain{}, \CvgObsAbstBola{} of
them true violations. Within the records it decides it detects
\CvgObsDetected{} and raises no false proof. On the \CvgJointVthreeDirect{} records
both methods decide, it detects \CvgDetectGap{} more violations than v3; the
remaining difference in pooled recall is coverage, not accuracy. The last row of
Table~\ref{tab:coverage} is the binary version that was implicit before: it
forces a negative verdict where the probe is unreadable, covers all
\CvgN{} records, detects the same \CvgObsDetected{}, and misses
\CvgForcedObsFN{}. Its selective risk of \CvgForcedObsRiskPct{}\% is entirely
missed violations, produced by the forcing rather than by the method; we report
it as a reference point, not as a comparator.

Selective risk is not a ranking either, and we do not use it as one.
\texttt{degrade\_all}, which alarms on nothing, has a selective risk of
\CvgDegradeRiskPct{}\% at \CvgDegradeCoveragePct{}\% coverage and detects no
violation at all; an ordering by risk alone would place it above v2. We therefore
report precision and selective risk with coverage attached and never compare them
across methods at unequal coverage. The stratified analysis
(Section~\ref{sec:stratum}) equalises the \emph{availability} of the oracle, not
the decision rate: v3 still declines \CvgVthreeCover{} readable records that show
a covering grant, so the two strata do not make coverage common, and we say so
wherever the comparison is made. A safe abstention is not a detection and we
never count it as one.

\paragraph{Cost} The enumeration issues \CvgReqTotal{} requests over \CvgN{}
records (median \CvgReqMedian{} per record, range \CvgReqMin--\CvgReqMax). The
oracle-consuming methods add \CvgProbeGets{} probe GETs over \CvgProbeCells{}
cell-operations: charged per record that is one request, and \CvgProbePerRecord{}
amortised as executed, because the probe is shared by the records of a cell.
\texttt{raw}, v1 and v2 issue none of their own---they read the log the crawl
already wrote. We do not report wall-clock time for this study: every method
replays the same recorded exchanges, so a difference would measure our replay
loop rather than the system under test. The study we did execute end to end is
the temporal loop of Section~\ref{sec:realsystem}: a median of
\CvgLoopGiteaTotal{}~s per transition on Gitea and \CvgLoopGitlabTotal{}~s on
GitLab, of which the mutation is \CvgLoopGiteaMut{}~s and \CvgLoopGitlabMut{}~s,
the recrawl \CvgLoopGiteaRecrawl{}~s and \CvgLoopGitlabRecrawl{}~s, and the
cross-credential probe sweep \CvgLoopGiteaProbe{}~s and \CvgLoopGitlabProbe{}~s.
Each of these is a separate median, so the parts need not sum to the total
exactly. The recrawl the mechanism needs is a small part of the loop; re-reading
every credential is what costs.

\input{coverage_table.tex}

\subsection{A published metric's resolution: what a pooled specificity cannot register}
\label{sec:resolution}

The published figure is a specificity, not a false-alarm share of all inputs.
\mbox{MST-wi} defines it as $\mathrm{TN}/(\mathrm{TN}+\mathrm{FP})$ over the
follow-up inputs $F$, and states that each reported MT failure was manually
verified~\cite{Chaleshtari2022Metamorphic}; the labels are therefore not the
crawler's own answer. With $\mathrm{FN}=0$ throughout that evaluation,
$\mathrm{TN}=F-\mathrm{TP}-\mathrm{FP}$ and the figure equals
$1-\mathrm{FP}/(F-\mathrm{TP})$: the denominator is the \emph{negative}
follow-up inputs, not every generated input. We use that form below, and we do
not attribute the published number to a circular label.

What such a figure cannot do is register a failure confined to a class, and the
two reasons are arithmetic and coverage rather than labelling. First, the figure
is pooled: a class contributing a share $w$ of the negative inputs can move it
by at most $w$, so a real class-uniform failure is diluted by unaffected
inputs. Second, the follow-up inputs are those the crawl could enumerate; a
class of legitimate access the crawl never reaches contributes no negative
inputs and therefore cannot enter the ratio at all. The second point is the
measurement face of Finding~1, and it is a statement about the coverage of the
input set, not about the correctness of the human labels. We cannot reproduce
the engine's exact input count, so we substitute no value for the published
figure; the argument below is about the form of the denominator only.

\paragraph{A bounded formal remark on the denominator}
Write $S=1-\mathit{FP}/N^{-}$, where $N^{-}=\mathrm{TN}+\mathrm{FP}$ is the
number of negative follow-up inputs the run generated, and define
$S_{\mathrm{scenario}} = 1 - N_{\mathrm{failed}} / N_{\mathrm{scenarios}}$,
where $N_{\mathrm{failed}}$ is the number of scenarios in which the failure is
uniform. $S_{\mathrm{scenario}}$ is an analytical construct for characterising
resolution; it is not proposed as a replacement metric. The remark concerns
resolution, not estimator consistency or power.

\begin{proposition}
\label{prop:resolution}
Fix $f \ge 1$ failed scenarios, each producing a fixed number $a \ge 1$ of
false-alarm inputs, and $s \ge 1$ successful scenarios each contributing $m$
additional alarm-free negative inputs. Then
$S = 1 - (f a) / (f a + s m)$, which tends to $1$ as $m \to \infty$, while
$S_{\mathrm{scenario}} = 1 - f/(f+s)$ is independent of $m$ and strictly less
than $1$.
\end{proposition}

\begin{proof}
$N^{-} = f a + s m$ and $\mathit{FP} = f a$; substitution gives the first
expression, which increases monotonically in $m$ to $1$, whereas the
scenario-level figure depends only on the fixed counts $f$ and $s$.
\end{proof}

\begin{corollary}
\label{cor:resolution}
No fixed input-level specificity threshold $t < 1$ can guarantee that a
scenario-uniform failure will push the metric below it: choosing
$m > (f a/s)\cdot t/(1-t)$ drives $S$ above $t$ while the $f$ failing
scenarios remain failing.
\end{corollary}

The finding is therefore about what the published figure can be asked to show,
not about its arithmetic: it is pooled over classes, and it is computed on the
inputs the crawl reached. Neither point disputes the reported value or its
labels; both limit what a high value certifies about a class the run did not
exercise.

\paragraph{Implementation robustness}
An independently written implementation of the same predicates was compared on a
subset fixed in advance; it bounds implementation accident only, and is reported
with the confirmations in Section~\ref{sec:secondimpl}.
